\subsection{Sparse projected addresses and disjoint sampling slots}

We now give a finite formulation of the address count.  Throughout this
subsection, $\mathbb P$ denotes the raw bilateral inventory-word law.
Additional randomness used to select finite witness lists may be included,
provided that it does not change the word's marginal law.  A subsequent
selection law $\mathbb Q_N$ will be required to satisfy
\begin{equation}\label{arm:address-tilt}
 \frac{\mathrm d\mathbb Q_N}{\mathrm d\mathbb P}
 \leq \exp(k_N),\qquad k_N=o(N).
\end{equation}
The independent grid displacement introduced below has the same conditional
law under both measures.

\begin{definition}[Projected address forest]\label{arm:address-forest}
Fix an integer $R\geq8$ and write $\lambda=\log R$.  Let $I$ be a
fixed interval of length $H$, where $H$ is a fixed constant multiple of the
base time unit $n$.  Let
\[
 J=\left\lfloor\frac{bN}{\lambda}\right\rfloor,
 \qquad \delta_j=H R^{-j},\quad 0\leq j\leq J.
\]
All retained time units are assumed to diverge.  Select a finite list of
markers $(t,\ell)$ with $t\in I$ and $0\leq\ell\leq J$.  The list,
including the exact positions, terminal levels, and all incidences, may
be a function of the entire word.  \emph{After this selection}, sample
$U$ uniformly on $[0,H)$, independently of everything already selected.
Use the nested half-open grids
\[
 \mathcal G_j(U)=
 \big\{[U+k\delta_j,U+(k+1)\delta_j):k\in\mathbb Z\big\}.
\]
A marker $(t,\ell)$ occupies the containing cell at each level $j\leq\ell$.
Let $V_j$ be the set of occupied cells, and put
\[
 m_j=|V_j|,\qquad T=\sum_{j=0}^Jm_j,
 \qquad A=\lambda T.
\]
The occupied cells, with the containing-cell parent relation, form the
\emph{projected address forest}.  Write $M$ for the total number of marker
entries, with multiplicity; entries terminating at an internal forest
vertex are retained as terminal labels there.

There are $h_0$ distinguished bases, each marked through level $J$.
At a probe level $j_p$, a witness list consists of finitely many chains
from these bases.  Every chain endpoint and intermediate position is
marked through $j_p$.  Every link is either a chord of one of the two
colors, or a time jump satisfying
\[
 |s-t|\leq K\delta_{j_p}.
\]
Here $K$ is a fixed finite list bound, enlarged to bound the number of
marker and edge entries per probe as well.  Thus each chain projects in
its entirety to every coarser level.  Exact endpoints within their
terminal cells will not be part of the address code.
\end{definition}

Only a bounded number of root cells can meet $I$; with the convention in
Definition~\ref{arm:address-forest}, there are at most two.  Their integer
names are determined by $I$ and $U$, so no unspecified macroscopic time
coordinate must be encoded.  All the sampling slots below are contained
in a fixed deterministic enlargement of $I$.  The bilateral word supplies
their chronological increments, including increments outside $I$ if a
slot meets an endpoint of that interval.

\begin{lemma}[Projected graph ranks]\label{arm:projected-rank}
At level $j$, retain all occupied boxes, and place the projected colored
chords between them, including isolated vertices.  Let $c_j$ be the
number of components of the union graph, let $b_j$ be the number of
occupied base boxes, and let $D_j$ count active jump entries whose two
endpoints lie in different level-$j$ boxes.  If $e_j$ is the sum of the
two monochrome graph ranks, then
\begin{equation}\label{arm:rank-component-bound}
 c_j\leq b_j+D_j,
 \qquad e_j\geq m_j-c_j\geq m_j-b_j-D_j.
\end{equation}
These counts are made before any sampling slots are discarded.
\end{lemma}
\begin{proof}
Include the projected jump edges temporarily.  Every occupied vertex
then has a chain to a base: a marker arises in a fully retained chain,
or is itself a base.  The enlarged graph therefore has at most $b_j$
components.  Removing the $D_j$ jump edges increases the component count
by at most $D_j$.  For the rank inequality, take a spanning forest of the
remaining union graph.  Its edges of either color form a forest in that
monochrome graph, so their number is bounded by that color's rank.
Summing the two bounds gives $m_j-c_j\leq e_j$.
\end{proof}

\begin{lemma}[Regular slots]\label{arm:regular-slots}
For every $\varepsilon>0$, one can choose a fixed sufficiently large
$R$ such that, conditional on any selected marker system, with
probability at least $7/8$ over $U$, all but
\begin{equation}\label{arm:regular-deletion}
 C_R M+\varepsilon T
\end{equation}
of its occupied vertex occurrences admit the following disjoint paired
slots.  At a retained vertex $v\in V_j$, let $t_v$ be the midpoint of
its cell and put $r_j=R^2\delta_j$.  The slots are
\begin{equation}\label{arm:paired-slots}
 S_v^-=[t_v-r_j,t_v-2r_j/R],
 \qquad
 S_v^+=[t_v+2r_j/R,t_v+r_j].
\end{equation}
All these intervals, for every retained vertex, level, and sign, are
pairwise disjoint.  Their locations are determined by the forest and
$U$.  The constant $C_R$ is independent of the positions, the number of
probe levels, and the finite exploration cutoff except through $M$.
\end{lemma}
\begin{proof}
Call a vertex critical if it is a root, has a number of children different
from one, or carries a terminal entry.  There are at most $M$ leaves and
at most $M$ terminating vertices.  For a finite forest,
\[
 \sum_{v:\,\deg^+(v)\geq2}(\deg^+(v)-1)
 =\#\{\text{leaves}\}-\#\{\text{roots}\},
\]
so there are at most $M$ branching vertices as well.  Thus at most $4M$
vertices are critical.  Fix an integer $g=4$ and discard all vertices at
forest distance at most $g$ from a critical vertex.  Since the degree is
at most $R+1$, this discards at most
\[
 C_RM,\qquad
 C_R=4\sum_{k=0}^g(R+1)^k.
\]

For a remaining vertex $v\in V_j$, its ancestor $Q_v$ at depth $j-g$
exists.  The intervening path is unary.  In particular, $v$ is the only
occupied depth-$j$ descendant of $Q_v$.  Additionally discard $v$ if
its cell is at distance less than $3r_j$ from $\partial Q_v$.
Call this a boundary deletion.  We first prove disjointness for the
vertices left by these two deletion rules.

Take distinct retained vertices $v\in V_j$ and $w\in V_k$, with
$k\geq j$.  If $k=j$, the cell of $w$ is outside $Q_v$, since $v$ is
its sole occupied depth-$j$ descendant.  Hence
$|t_w-t_v|\geq3r_j$, whereas the sum of outer slot radii is $2r_j$.
If $k>j$ and $w$ descends from $v$, then
\[
 |t_w-t_v|\leq\delta_j,\qquad r_k\leq r_j/R,
 \qquad
 \delta_j+r_j/R
 =r_j(R^{-2}+R^{-1})<2r_j/R.
\]
Both slots at $w$ therefore lie strictly inside the inner gap of the
slots at $v$.  If $k>j$ and $w$ does not descend from $v$, its occupied
ancestor at depth $j$ lies outside $Q_v$.  Thus $t_w$ is outside $Q_v$,
$|t_w-t_v|\geq3r_j$, and the sum of outer radii is at most
$r_j+r_j/R<2r_j$.  This proves disjointness in every case.  The two
slots at a single vertex are disjoint by their definition.

It remains to count boundary deletions without imposing any distribution
on the marker positions.  Condition on all these positions and terminal
levels.  At each depth $j$, let $s_j$ be the occupancy count for a fixed
auxiliary grid of width $\delta_j$, with no random displacement, and
let $S=\sum_j s_j$.  Each cell of either of two grids of the same width
meets at most two cells of the other grid.  It follows, for every $U$,
that
\begin{equation}\label{arm:grid-comparability}
 S/2\leq T(U)\leq2S.
\end{equation}
Choose one active marker in each occupied auxiliary cell.  If an
occupied primary cell at depth $j\geq g$ is boundary-deleted, a chosen
representative in an auxiliary cell meeting it lies within
$3r_j+3\delta_j$ of a depth-$(j-g)$ boundary.  The displacement modulo
$\delta_{j-g}$ is uniform.  The probability of this event for a fixed
representative is at most
\[
 \frac{2(3r_j+3\delta_j)}{\delta_{j-g}}
 \leq12R^{2-g}.
\]
An auxiliary cell meets at most two primary cells.  If $B(U)$ is the
number of boundary deletions, linearity of expectation gives
\[
 \mathbb E_U B(U)\leq24R^{2-g}S.
\]
By Equation~\eqref{arm:grid-comparability} and Markov's inequality,
\begin{equation}\label{arm:boundary-probability}
 \mathbb P_U\{B(U)>\varepsilon T(U)\}
 \leq\frac{48R^{2-g}}{\varepsilon}.
\end{equation}
For $g=4$, choose $R$ large enough that this is at most $1/8$.
Together with the structural deletions this proves
Equation~\eqref{arm:regular-deletion}.
\end{proof}

Figure~\ref{fig:regular-slots} illustrates the nested-slot separation and
the padding inside the unary ancestor.
\begin{figure}[tbp]
\centering
\begin{tikzpicture}[
  x=1cm,y=1cm,
  font=\small,
  >=Latex,
  coarse/.style={draw=blue!65!black,fill=blue!13,line width=.8pt},
  finer/.style={draw=teal!70!black,fill=teal!20,line width=.8pt},
  guide/.style={draw=black!35,densely dashed,line width=.45pt},
  dimension/.style={<->,draw=black!65,line width=.5pt},
  every node/.style={inner sep=2pt}
]
\node[anchor=west,font=\footnotesize] at (.35,4.18)
  {Cell nesting (magnified, schematic)};
\draw[draw=black!55,rounded corners=1.5pt] (.35,2.88) rectangle (14.05,3.88);
\node[anchor=west] at (.55,3.65) {$Q_v$};
\node[anchor=east,font=\footnotesize,text=black!65] at (13.85,3.65)
  {unary ancestor at depth $j-4$};
\draw[coarse] (5.85,3.04) rectangle (8.55,3.69);
\node[anchor=west,text=blue!65!black] at (6.0,3.37) {$v$};
\draw[finer] (7.26,3.14) rectangle (7.94,3.59);
\node[text=teal!70!black] at (7.60,3.37) {$w$};
\draw[dimension] (.55,3.13) -- (5.65,3.13)
  node[midway,above,font=\footnotesize] {$\ge 3r_j$};
\draw[dimension] (8.75,3.13) -- (13.85,3.13)
  node[midway,above,font=\footnotesize] {$\ge 3r_j$};

\fill[black!4] (5.75,-.70) rectangle (8.65,2.20);
\draw[guide] (5.75,-.70) -- (5.75,2.20);
\draw[guide] (8.65,-.70) -- (8.65,2.20);
\draw[->,draw=black!60] (.35,-.47) -- (14.05,-.47);
\node[anchor=east,font=\footnotesize] at (14.05,-.76) {time};

\draw[coarse] (1.4,1.03) rectangle (5.75,1.45);
\draw[coarse] (8.65,1.03) rectangle (13.0,1.45);
\node[text=blue!65!black] at (3.575,1.24) {$S_v^-$};
\node[text=blue!65!black] at (10.825,1.24) {$S_v^+$};
\node[anchor=south] at (1.4,1.62) {$t_v-r_j$};
\node[anchor=south] at (5.75,1.62) {$t_v-2r_j/R$};
\node[anchor=south] at (8.65,1.62) {$t_v+2r_j/R$};
\node[anchor=south] at (13.0,1.62) {$t_v+r_j$};
\draw[draw=blue!65!black,line width=.65pt] (7.20,.85) -- (7.20,1.53);
\fill[blue!65!black] (7.20,1.24) circle (1.6pt);
\node[anchor=south,text=blue!65!black] at (7.20,1.60) {$t_v$};

\draw[finer] (6.505,.05) rectangle (7.04875,.47);
\draw[finer] (7.41125,.05) rectangle (7.955,.47);
\node[anchor=south,text=teal!70!black] at (6.77,.52) {$S_w^-$};
\node[anchor=south,text=teal!70!black] at (7.69,.52) {$S_w^+$};
\draw[draw=teal!70!black,line width=.65pt] (7.23,-.14) -- (7.23,.53);
\fill[teal!70!black] (7.23,.26) circle (1.6pt);
\draw[draw=teal!70!black,line width=.45pt] (7.23,-.14) -- (7.67,-.30);
\node[anchor=west,text=teal!70!black] at (7.67,-.30) {$t_w$};
\node[anchor=west,font=\footnotesize] at (.55,.26)
  {finer descendant, $k>j$};
\draw[dimension] (5.75,-1.02) -- (8.65,-1.02);
\node[anchor=north,font=\footnotesize] at (7.20,-1.08)
  {coarse central gap};
\end{tikzpicture}
\caption[Disjoint regular sampling slots]{%
Regular sampling slots from Lemma~\ref{arm:regular-slots}.
The inset shows cell containment; the two slot rows share one time axis
and are vertically separated for visibility.  Here $r_j=R^2\delta_j$,
the unary ancestor $Q_v$ lies at depth $j-4$, $v$ is its sole occupied
depth-$j$ descendant, and $\operatorname{dist}(v,\partial Q_v)\ge3r_j$.
For a finer descendant $w\subset v$ at depth $k>j$,
$|t_w-t_v|+r_k\le\delta_j+r_j/R
=r_j(R^{-2}+R^{-1})<2r_j/R$;
thus both finer slots lie strictly inside the coarse central gap,
although they need not lie inside the cell $v$.
For an incomparable retained node at depth $k\ge j$,
$|t_w-t_v|\ge3r_j>r_j+r_k$.
These inequalities give disjointness in both cases.}
\label{fig:regular-slots}
\end{figure}


\begin{remark}[Other fixed test bands]\label{arm:other-test-bands}
The same construction applies to a fixed finite family of test bands
contained in $a_*r\leq |t-t_v|\leq A_*r$, where
$0<a_*<A_*<\infty$.  First refine the location boxes until their width
is at most $a_*r/4$, and sample generations far enough apart that the
next outer radius is at most $a_*r/4$.  This uses a fixed larger grid
ratio or a fixed residue class of generations; an independent uniform
choice of the residue preserves uniform sampling of levels.  Increase
the boundary padding beyond $2A_*r$ and then increase the fixed ancestor
depth $g$ if necessary.  A finer descendant's bands are now inside the
coarser inner gap, while an incomparable descendant lies outside the
padded ancestor.  The preceding disjointness proof applies verbatim to
these two explicit inequalities.  The structural and boundary counts
have the same form, with constants depending on the fixed bands.
For integer word times, rounding slot boundaries inward preserves
disjointness and changes normalized lengths by a vanishing amount.
\end{remark}

\begin{lemma}[Integrated connectivity errors]\label{arm:jump-errors}
For each fixed $K,R$, a jump entry at terminal level $\ell$ contributes
at most $C_{K,R}<\infty$ to the expected sum of its separated coarse
projections.  Consequently, if $M_{\mathrm{jump}}$ is the number of jump
entries,
\begin{equation}\label{arm:jump-sum}
 \mathbb E_U\sum_{j=0}^JD_j
 \leq C_{K,R}M_{\mathrm{jump}}.
\end{equation}
If the bases lie in an interval of length at most $C H e^{-N}$ and
$j_N=\lfloor N/\lambda\rfloor$, then also
\begin{equation}\label{arm:base-merger-error}
 \mathbb E_U\sum_{j=0}^{j_N}(b_j-1)\leq C'_{C,R}.
\end{equation}
For a sparse family with $P=O(N/L)$ probes, $L\longrightarrow\infty$,
and fixed list bound $K$, there is a deterministic $d_N=o(N)$ such that,
uniformly in the selected witness data,
\begin{equation}\label{arm:good-connectivity}
 \mathbb P_U\left\{
 \sum_{j=0}^JD_j+
 \sum_{j=0}^{j_N}(b_j-1)>d_N\right\}=o(1).
\end{equation}
Moreover $C_RM=o(N)$.  Together with
Lemma~\ref{arm:regular-slots}, these assertions hold simultaneously on
a set of displacements of conditional probability at least $1/2$ for
all sufficiently large $N$.
\end{lemma}
\begin{proof}
Two fixed times at distance $d$ lie in different uniformly displaced
cells of width $\delta_j$ with probability at most
$\min\{1,d/\delta_j\}$.  Therefore a jump of length at most
$K\delta_\ell$ has expected total separated projections at most
\[
 \sum_{j=0}^\ell\min\{1,K R^{j-\ell}\}
 \leq \sum_{a=0}^\infty\min\{1,K R^{-a}\}
 =:C_{K,R}<\infty.
\]
Linearity proves Equation~\eqref{arm:jump-sum}.  For
Equation~\eqref{arm:base-merger-error}, list the bases in increasing
order and join consecutive bases by artificial edges for counting only.
At every depth, $b_j-1$ is bounded by the number of these edges that
cross a cell boundary.  Their total Euclidean length is at most
$CHe^{-N}$.  Hence
\[
 \mathbb E_U(b_j-1)\leq C e^{-N}R^j,
 \qquad
 \sum_{j=0}^{j_N}C e^{-N}R^j\leq \frac{CR}{R-1}.
\]
For fixed $K$, both $M$ and $M_{\mathrm{jump}}$ are bounded by a fixed
multiple of $P+h_0$.  Choose $d_N=o(N)$ so slowly that
$(P+h_0+1)/d_N\longrightarrow0$.  Markov's inequality gives
Equation~\eqref{arm:good-connectivity}; the bound $C_RM=o(N)$ follows
from $P=o(N)$.  The last assertion follows by intersecting this event
with the event of Lemma~\ref{arm:regular-slots}.
\end{proof}

\begin{lemma}[Stability under an independent grid displacement]
\label{arm:grid-partition-stability}
Fix a marker system with $M$ entries before sampling its displacement.
At level $j$, partition its active entries without using a grid: first
group entries at the same time, then join consecutive distinct active
positions whenever their gap is strictly less than $\delta_j$.
Denote the resulting partition by $\mathcal P_j^*$, and let
$\mathcal P_j(U)$ be the partition into occupied cells of the displaced
grid.  Put
\[
 \Delta(U)=\#\{0\le j\le J:\mathcal P_j(U)\ne\mathcal P_j^*\}.
\]
For an independent uniform displacement,
\begin{equation}\label{arm:grid-partition-error}
 \mathbb E_U[\Delta(U)\mid\text{marker system}]
 \le \frac{2R}{R-1}M.
\end{equation}
At every level where the partitions agree, the entire projected
colored graph, including its base labels, is the same up to cell names.
In particular, if $M\le M_N=o(N)$, two independent uniform
displacements give the same projected graphs outside $o(N)$ levels
with probability tending to one, uniformly over the selected marker
system.
\end{lemma}

\begin{proof}
If $M=0$, all partitions are empty and the assertions are immediate.
Keep the labels and incidences of coincident entries when grouping
their positions.  As $j$ increases, entries and hence distinct active
positions can only disappear.  Initially there are at most $M-1$
consecutive-position pairs; deleting one position creates at most one
new pair.  Thus at most $2M-1$ distinct consecutive pairs occur over
all levels.  Deleting all positions with a common terminal level one
at a time only enlarges this bound.

Every grid partition refines $\mathcal P_j^*$: positions in the same
half-open cell have total span strictly less than $\delta_j$, and so
do all their successive gaps.  The partitions can differ only if a
consecutive active pair with gap $d<\delta_j$ is split by a cell
boundary.  For a fixed pair this has probability $d/\delta_j$, since
$U$ modulo $\delta_j$ is uniform.  For each positive gap,
\[
 \sum_{j:\,\delta_j>d}\frac d{\delta_j}
 \le \frac R{R-1}.
\]
The same bound holds after restricting to the levels at which the
pair is consecutive and active.  A union bound over pairs and levels
proves Equation~\eqref{arm:grid-partition-error}.

All active edge incidences and base labels were fixed before the grid,
so equal marker partitions give identical projected graph data.
Choose a deterministic $q_N=o(N)$ with $M_N/q_N\to0$.
Markov's inequality bounds the conditional probability of
$\Delta(U)>q_N$ by $2RM_N/((R-1)q_N)=o(1)$.  Apply this to each of
two independent displacements.  Outside that exceptional event they
both agree with $\mathcal P_j^*$ except at at most $2q_N$ levels.
The estimate is conditional on arbitrary marker data, so it also
holds under any selection law sampled before the displacements.
\end{proof}

\begin{remark}[Why the order of selection matters]
The witnesses in Definition~\ref{arm:address-forest} must precede $U$.
For example, take two terminal positions $U-\eta$ and $U+\eta$ with
$\eta\ll\delta_J$.  They occupy distinct unary branches on opposite
sides of the persistent boundary $U$ at all $J+1$ levels, but the
same-level paired slots overlap.  At least half the $2(J+1)$ nodes
must be deleted, although there are only two terminal entries.  This
contradicts a bound $C_RM+\varepsilon T$ with $\varepsilon<1/2$ for
large $J$.  Our conditional estimates avoid this obstruction.  We keep
$U$ independent in the probability argument, rather than treating a
sample-dependent favorable displacement as free fixed information.
\end{remark}

\Needspace{8\baselineskip}
\subsection{Counting the address codes}

\begin{lemma}[Address code count]\label{arm:address-codes}
Fix $U$ and the distinct probe depths $j_1,\ldots,j_P$.  Include the
base entries in the probe list at depth $J$, adding that depth if needed,
and enlarge $K$ to cover those entries.  The resulting probe depths
remain distinct.  Suppose each list has at most $K$ marker and edge entries.  There is a coding of the
projected forests and their list incidences with the following bound,
uniformly in the total occupancy $T$:
\begin{align}
 \log\#\{\text{codes of occupancy }T\}
 &\leq(\lambda+C_0)T
   +KP\log(1+T/P)
   +CPK\log(K+1)+O(1).\label{arm:address-code-count}
\end{align}
The constant $C_0$ does not depend on $K$.  A family of additional
annotations with at most $C_a^{m_j}$ choices at depth $j$, for a fixed
$C_a$, only increases $C_0$ by $\log C_a$.
In particular, for fixed $K$ and $P=o(N)$,
\begin{equation}\label{arm:address-code-leading}
 \log\#\{\text{codes of occupancy }T\}
 \leq(\lambda+C_0)T+o(T+N),
\end{equation}
uniformly in $T$.  A sufficient condition for a growing bound $K=K_N$
and probe stride $L$ is
\begin{equation}\label{arm:growing-list-condition}
 \frac{K_N\log(K_NL)}{L}\longrightarrow0.
\end{equation}
The codes can also be assigned subprobability weights at an additional
cost $2\log(T+1)+O(1)$.
\end{lemma}
\begin{proof}
An ordered forest with $T$ vertices and at most two roots has at most
$C4^T$ possible shapes.  One way to see this is to attach a new root
above its roots and use the depth-first parenthesis word of the
resulting ordered tree; its length is $2(T+1)$, and a balanced
parenthesis word determines the tree.  Each nonroot cell has one of
$R$ child offsets in its parent, so the shape has at most $R^T$
geometric assignments.  The occupied root cells have a bounded presence
mask among the cells meeting $I$, whose names are already fixed by
$I,U$.  This proves the bound $C(4R)^T$ for the forest before labels.
Its vertices' depths follow from its parent relation.

At probe depth $j_p$, name a position by its ordered index among the
$m_{j_p}$ occupied cells; use one extra index for an absent list entry.
Thus the position entries have at most $(1+m_{j_p})^K$ possibilities.
Repeated indices record coincident positions.  Names in different
probe lists do not require global forest indices: each local index
specifies a forest vertex, whose ancestors are already determined by
the forest.  A correspondence to a distinguished base uses that base's
fixed name, and its depth-$J$ address determines its coarser addresses.
The list's chain order, colors, endpoint incidences, and base names have
at most $\exp(CK\log(K+1))$ possibilities.  The count includes absent
entries and shorter lists.  It never specifies an exact endpoint
within its terminal cell.

Multiply these choices and use the distinctness of probe depths to
obtain $\sum_p m_{j_p}\leq T$.  Concavity gives
\[
 \sum_{p=1}^P\log(1+m_{j_p})
 \leq P\log(1+T/P),
\]
which proves Equation~\eqref{arm:address-code-count}.  The same estimate
bounds the maximum number of list decorations over all forest shapes,
so variation of the individual $m_j$ between shapes causes no extra
factor.

For $0<\eta\leq K$ and $z\geq0$,
\[
 K\log(1+z)\leq\eta z+K\log(K/\eta).
\]
Indeed the maximum of $\log(1+z)-(\eta/K)z$ occurs at
$z=K/\eta-1$ and is at most $\log(K/\eta)$.  Hence
\begin{equation}\label{arm:address-code-formbound}
 \log\#\{\text{codes of occupancy }T\}
 \leq(\lambda+C_0+\eta)T
   +KP\log(K/\eta)+CPK\log(K+1)+O(1).
\end{equation}
For fixed $K$, let $\eta$ tend to zero sufficiently slowly while
$P=o(N)$.  This yields Equation~\eqref{arm:address-code-leading}
uniformly in $T$.  Under Equation~\eqref{arm:growing-list-condition},
for example $\eta_N=1/\log L$ for large $L$ makes both remainders in
Equation~\eqref{arm:address-code-formbound} equal to $o(N)$.
No condition $L\gg\log N$ is used.  If the probe depths themselves
must be encoded among the $J+1$ levels, their extra cost is at most
$\log\binom{J+1}{P}\leq P\log(e(J+1)/P)=o(N)$ for
$P=O(N/L)$; alternatively their independent mesh displacement can be
conditioned on.

Finally give $T$ a weight proportional to $(T+1)^{-2}$ and distribute
that weight uniformly among its finite enumeration of codes.  The
weights sum to at most one, and the negative logarithm of each weight
is the code count bound plus $2\log(T+1)+O(1)$.
\end{proof}

We call a displacement \emph{favorable} if it satisfies the boundary
bound of Lemma~\ref{arm:regular-slots} and the connectivity bound of
Lemma~\ref{arm:jump-errors}.  For fixed $R,K$, its conditional
probability given the selected witness system is at least $1/2$ for
large $N$.  At a favorable displacement the structural deletions are
$o(N)$ and
\begin{equation}\label{arm:total-rank-lower}
 E:=\sum_{j=0}^Je_j
 \geq T-h_0(J+1)-d_N.
\end{equation}
For earlier scales, Equation~\eqref{arm:base-merger-error} also gives
the stronger replacement of $h_0$ by one, up to an integrated $o(N)$
error.  The simpler bound in Equation~\eqref{arm:total-rank-lower}
suffices for the following occupancy estimate.
